\section{What on Earth is Sequence Transduction?}

The abstract opens with: \emph{``The dominant sequence transduction models are based on complex recurrent or convolutional neural networks that include an encoder and a decoder.''}

There is a lot going on in that sentence, but the first thing to deal with is this phrase \emph{sequence transduction}, because if you don't know what it means, the rest of the paper is going to feel like reading a novel where you missed the first chapter.

So.
A \textbf{sequence} is just an ordered list of things.
A sentence is a sequence of words.
A song is a sequence of notes.
A conversation is a sequence of utterances.

\textbf{Transduction} is a fancy word for conversion.
A microphone transduces sound into electrical signals.
Your retina transduces light into nerve impulses.

Put them together and \textbf{sequence transduction} is: you take a sequence in, you get a different sequence out.
That is it.
The input sequence and the output sequence need not be the same length, or in the same language, or even the same kind of thing.

The canonical example is machine translation.
English in, French out.
But once you see the pattern, it is everywhere:

\begin{itemize}
    \item Speech recognition: audio waveform in, text out.
    \item Summarization: a long document in, a short one out.
    \item Chatbots: your message in, a reply out.
\end{itemize}

Now, the reason I belabor this point is that the paper says \emph{dominant} sequence transduction models.
Dominant.
Not ``all'' or ``the only.''
The authors are being very specific: they are describing the state of the art in 2017, and the state of the art had converged on a particular family of approaches.

Which raises the obvious question: what were those approaches, and why had everyone converged on them?
